\section{Direction Features and Explicit Joint Measurements}\label{app:coordinate-control}
\subsection{Purpose and Relationship to Shin}
The original motivation was to test descriptions such as bent, extended, spread and touching. Continuous measurements are a comparison with these categories. Shin's 630 direction cosines already encode all 210 landmark-pair directions. For $u=(a-b)/\|a-b\|$ and $v=(c-b)/\|c-b\|$, the joint angle is $\arccos(u^\top v)$: it is obtainable from those descriptors with appropriate pair signs. Thresholding an angle into a category likewise supplies no independent sensor information.

Making a nonlinear relation explicit may help a classifier. The original additions also use estimated world landmarks, while Shin uses normalized image landmarks. The control separates coordinate source from joint representation; it does not assume that Shin omitted finger-bending information.

\subsection{Matched Design}
The same 60,251 accepted images and five frozen filename-block folds are retained, without new detection, augmentation or tuning. The 15 measurements (14 interior angles plus signed thumb-MCP deviation) use Shin's exact normalized image $x,y,z$ coordinates without aspect correction. These are descriptors, not calibrated anatomical angles.

The reference has 190 distances and 630 directions (820 inputs). Replacement substitutes 15 joint measurements for directions (205 inputs); addition retains the reference (835). The earlier 835-input world-joint model provides context. Historical predictions are aligned by ID, label and fold, with source hashes checked. Imputation and standardization use training rows only. Fixed classifier settings match the original profiles; RBF \texttt{gamma=scale} adapts to each representation. These recipes are not equally optimized upper bounds.

\begin{table}[h]
\centering\small
\setlength{\tabcolsep}{4pt}
\caption{Coordinate control on matched held-out predictions. Fixed/lost/net refer to each learner's Shin reference. The world rows reproduce the original comparison; they are not new runs.}
\label{tab:coordinate-control}
\begin{tabular}{llrrrrrr}\toprule
Learner & Representation & Inputs & Acc. (\%) & Macro-F1 & Fixed & Lost & Net\\\midrule
\tableinput direction_joint_rows.tex
\bottomrule\end{tabular}
\end{table}

\subsection{Interpretation and Evidence}
Adding same-coordinate joints corrects 89 SVM errors and loses 85 correct predictions: net +4 (+0.0066 percentage points). Three folds improve and two worsen. The earlier world-joint result corrects 143 and loses 76: net +67 (+0.1112 points). Neither establishes new-person or webcam performance. LightGBM loses 34 correct predictions with same-coordinate joints and 48 with world joints.

Replacing directions loses a net 1,883 SVM and 1,901 LightGBM correct predictions, worsening every fold. For SVM, J-to-I errors increase from 28 to 531, I-to-J from 5 to 332, and H-to-U from zero to 230. Replacement removes axis-relative orientation as well as reducing feature count; this may discard useful letter distinctions. It is not evidence of faulty angle calculation, nor a controlled test of only rotation robustness.

The same-coordinate addition gives no convincing overall benefit under these recipes. The larger world-coordinate difference cannot be attributed to the bending formula alone; coordinate estimates, scaling and model geometry are possible contributors, not independently established causes. The full category/workbook package remains a separate comparison; this control does not isolate category-only predictive value.

The supplement \texttt{results/direction\_joint\_01/} contains per-letter metrics, confusion matrices, fold scores and per-image predictions. \texttt{research.paper\_core\_assets} generates paper rows after reconciling paired counts and historical scores. The implementation, runner and verifier are listed with commands in \texttt{REPRODUCE.md}. Computational checks do not establish independent scientific validation.
